% Generated by tools/make_interview_tracks.py. Do not edit.
\section*{Quels chapitres travailler}
Chaque case compte les questions d'un chapitre étiquetées pour ce métier ou ce type de firme. QR désigne le chercheur quantitatif (quant researcher), QT le trader quantitatif, DS le data scientist, QD le développeur quantitatif, Dev le développeur logiciel, LLE l'ingénieur basse latence, Risque le quant de risque~; TM le teneur de marché, Cr la firme crypto, Bq la banque, HF le hedge fund. Les épreuves et les exercices de vitesse sans étiquette sont pour tout le monde.
\par\medskip\noindent{\footnotesize\setlength{\tabcolsep}{2.5pt}%
\begin{tabular}{@{}>{\raggedright\arraybackslash}p{6.2cm}rrrrrrrrrrrrr@{}}
\hline
Chapitre & QR & QT & Trader & Analyste & DS & QD & Dev & LLE & Risque & TM & Cr & Bq & HF \\
\hline
\hyperref[ch:iv1:mental-technique]{Technique mentale} & 19 & 40 & 34 & 12 & 1 & 5 & 4 & 1 & 15 & 37 & 28 & 33 & 35 \\
\hyperref[ch:iv1:sequences-patterns-and-aptitude]{Suites, régularités et aptitudes} & 4 & 30 & 29 & 25 & 3 & 7 & 6 &  & 6 & 29 & 9 & 27 & 13 \\
\hyperref[ch:iv1:estimation]{Estimation} & 8 & 21 & 17 & 13 & 4 & 9 & 6 & 6 & 8 & 22 & 13 & 16 & 24 \\
\hyperref[ch:iv1:number-riddles]{Énigmes sur les nombres} & 24 & 25 & 17 & 1 &  & 11 & 5 &  &  & 25 & 14 & 14 & 25 \\
\hyperref[ch:iv1:counting-and-symmetry]{Comptage et symétrie} & 40 & 40 & 29 & 7 & 2 & 13 & 1 &  & 6 & 39 & 28 & 31 & 40 \\
\hyperref[ch:iv1:conditional-probability-and-bayes]{Probabilité conditionnelle et Bayes} & 38 & 40 & 24 & 14 & 16 &  &  &  & 18 & 32 & 17 & 24 & 39 \\
\hyperref[ch:iv1:expectation-and-indicators]{Espérance et variables indicatrices} & 25 & 23 & 16 & 3 & 6 & 4 & 4 & 1 & 7 & 25 & 20 & 17 & 25 \\
\hyperref[ch:iv1:waiting-times-and-markov-chains]{Temps d'attente et chaînes de Markov} & 25 & 25 & 15 & 3 & 2 & 1 &  &  & 6 & 22 & 13 & 14 & 25 \\
\hyperref[ch:iv1:random-walks-and-martingales]{Marches aléatoires et martingales} & 22 & 25 & 15 & 4 & 1 & 2 &  &  & 14 & 24 & 12 & 13 & 25 \\
\hyperref[ch:iv1:continuous-and-geometric-probability]{Probabilités continues et géométriques} & 25 & 25 & 16 & 3 & 9 & 5 & 1 & 1 & 3 & 24 & 17 & 18 & 25 \\
\hyperref[ch:iv1:order-statistics]{Statistiques d'ordre} & 15 & 15 & 3 & 4 & 10 & 2 &  & 1 & 4 & 14 & 11 & 10 & 14 \\
\hyperref[ch:iv1:optimal-stopping]{Arrêt optimal} & 20 & 20 & 17 & 5 & 5 &  &  &  &  & 15 & 11 & 11 & 19 \\
\hyperref[ch:iv1:games]{Jeux} & 20 & 20 & 14 & 6 &  & 4 &  &  &  & 20 & 14 & 13 & 19 \\
\hyperref[ch:iv1:logic-riddles]{Énigmes de logique} & 35 & 35 & 18 & 19 & 1 & 14 & 3 &  &  & 33 & 30 & 30 & 35 \\
\hyperref[ch:iv1:statistics]{Statistiques} & 40 & 40 & 5 & 16 & 35 & 1 &  & 1 & 10 & 28 & 24 & 25 & 40 \\
\hyperref[ch:iv1:regression]{Régression} & 30 & 22 & 2 & 4 & 25 & 2 &  &  & 16 & 24 & 20 & 25 & 30 \\
\hyperref[ch:iv1:linear-algebra]{Algèbre linéaire} & 25 & 22 &  &  & 16 & 15 & 1 &  & 10 & 14 & 14 & 16 & 25 \\
\hyperref[ch:iv1:calculus-and-optimisation]{Calcul et optimisation} & 20 & 16 & 3 & 4 & 8 & 9 & 4 & 1 & 1 & 15 & 14 & 14 & 20 \\
\hyperref[ch:iv1:machine-learning]{Apprentissage automatique} & 45 & 12 &  & 5 & 45 & 10 &  &  & 2 & 32 & 31 & 31 & 45 \\
\hyperref[ch:iv1:stochastic-calculus]{Calcul stochastique} & 30 & 30 &  &  &  & 24 &  &  & 18 & 8 & 8 & 23 & 30 \\
\hyperref[ch:iv1:time-series]{Séries temporelles} & 20 & 20 &  &  & 15 &  &  &  & 7 & 9 & 6 & 12 & 20 \\
\hline
\end{tabular}}
